\documentclass[11pt, a4paper]{article}

\usepackage[utf8]{inputenc}
\usepackage[T1]{fontenc}
\usepackage{amsmath, amssymb, amsthm, mathtools}
\usepackage{bm}
\usepackage{booktabs}
\usepackage{graphicx}
\usepackage{hyperref}
\usepackage[margin=2.5cm]{geometry}
\usepackage{enumitem}
\usepackage{xcolor}
\usepackage{tcolorbox}
\usepackage{float}

% Theorem environments
\theoremstyle{definition}
\newtheorem{definition}{Definition}[section]
\newtheorem{example}{Example}[section]
\theoremstyle{plain}
\newtheorem{theorem}{Theorem}[section]
\newtheorem{proposition}{Proposition}[section]
\newtheorem{corollary}{Corollary}[section]
\theoremstyle{remark}
\newtheorem{remark}{Remark}[section]

% Colored boxes
\newtcolorbox{keyresult}[1][]{
    colback=blue!5!white,
    colframe=blue!60!black,
    fonttitle=\bfseries,
    title={Key Insight},
    #1
}
\newtcolorbox{warningbox}[1][]{
    colback=red!5!white,
    colframe=red!60!black,
    fonttitle=\bfseries,
    title={Warning},
    #1
}
\newtcolorbox{refbox}[1][]{
    colback=gray!5!white,
    colframe=gray!50!black,
    fonttitle=\bfseries,
    title={Go Deeper},
    #1
}
\newtcolorbox{keyidea}[1][]{
    colback=green!5!white,
    colframe=green!50!black,
    fonttitle=\bfseries,
    title={Core Idea},
    #1
}
\newtcolorbox{prereq}[1][]{
    colback=orange!5!white,
    colframe=orange!60!black,
    fonttitle=\bfseries,
    title={Read Before Coding},
    #1
}
\newtcolorbox{objective}[1][]{
    colback=blue!5!white,
    colframe=blue!60!black,
    fonttitle=\bfseries,
    title={What You Will Build},
    #1
}

% Custom commands
\newcommand{\VaR}{\operatorname{VaR}}
\newcommand{\CoVaR}{\operatorname{CoVaR}}
\newcommand{\SRISK}{\operatorname{SRISK}}
\newcommand{\EV}{\mathbb{E}}
\newcommand{\PV}{\mathbb{P}}
\newcommand{\RV}{\mathbb{R}}
\newcommand{\ind}{\boldsymbol{1}}
\newcommand{\GCN}{\operatorname{GCN}}
\newcommand{\softmax}{\operatorname{softmax}}
\newcommand{\sigmoid}{\operatorname{sigmoid}}
\newcommand{\ReLU}{\operatorname{ReLU}}
\newcommand{\diag}{\operatorname{diag}}
\newcommand{\DebtRank}{\operatorname{DebtRank}}

\title{%
    \textbf{Graph Neural Networks for Credit Contagion} \\[0.5em]
    \large Modeling Systemic Risk Through Financial Networks
}
\author{Andr\'es Gonz\'alez Ortega}
\date{March 2026}

\begin{document}
\maketitle

\begin{abstract}
An applied treatment of Graph Neural Networks (GNNs) for modeling credit contagion
in financial networks.
Traditional credit risk models treat defaults as independent events or assume
dependence arises only from common macroeconomic factors.
The 2008 financial crisis demonstrated that interconnectedness itself is a source
of systemic risk: defaults cascade through networks of financial obligations.
We develop the framework from first principles: why network structure matters for
credit risk; how financial networks are represented as directed, weighted graphs;
the mechanisms of direct, fire-sale, and funding contagion; how GNNs learn
node-level representations by aggregating neighborhood information; Graph Convolutional
Networks (GCNs) for default prediction; the Eisenberg--Noe clearing mechanism for
generating cascade labels; and systemic risk metrics including DebtRank, CoVaR,
and SRISK\@.
A worked example runs throughout: a 50-node interbank network where GCN layers
propagate risk signals across the graph to predict which institutions default
following an exogenous shock.
We assume familiarity with neural networks, basic graph theory, and credit risk
fundamentals.
\end{abstract}

\tableofcontents
\newpage

%% ============================================================================
\section{Why Networks Matter for Credit Risk}
\label{sec:why-networks}
%% ============================================================================

Consider a stylized banking system with three banks: $A$, $B$, and $C$.
Bank~$A$ has lent \$500 million to Bank~$B$, and Bank~$B$ has lent \$400 million
to Bank~$C$.
An exogenous shock wipes out Bank~$C$'s assets.
Bank~$C$ defaults.
Bank~$B$ now loses \$400 million --- if this exceeds $B$'s capital buffer, $B$
defaults too.
Bank~$A$ then loses its \$500 million exposure to $B$, and the cascade continues.

This is \textbf{credit contagion}: the propagation of financial distress through a
network of contractual obligations.

\subsection{The Independence Assumption and Its Failure}

Traditional credit risk models rest on assumptions about how defaults are correlated:

\begin{itemize}
    \item \textbf{Merton (1974)}: each firm defaults when its asset value falls below
          its liabilities.
          Correlations arise from shared exposure to systematic factors (market returns).
          But the model treats each firm as an isolated entity --- it does not account
          for the contractual links between them.
    \item \textbf{CreditRisk+ (Credit Suisse, 1997)}: defaults are driven by
          independent Poisson processes, modulated by common risk factors.
          Dependencies come from shared factors, not from direct financial linkages.
    \item \textbf{Gaussian copula models}: a correlation parameter $\rho$ controls
          joint default probabilities, but the correlation is uniform and does not
          reflect the \emph{topology} of who owes whom.
\end{itemize}

These models capture the idea that a recession can cause many defaults simultaneously.
What they miss is that defaults cause other defaults \emph{mechanically}: through the
network of exposures.

\begin{warningbox}
The 2008 financial crisis provided a devastating empirical demonstration.
Lehman Brothers filed for bankruptcy on 15 September 2008 with approximately
930{,}000 derivative contracts outstanding.
Its default did not just cause losses at its direct counterparties --- it triggered
a chain reaction.
Counterparties that lost money on Lehman exposures became weaker, their creditors
became nervous, liquidity evaporated, and asset prices collapsed.
The total systemic losses far exceeded the direct losses on Lehman exposures.
\textbf{Interconnectedness amplified a single default into a global crisis.}
\end{warningbox}

\subsection{What Networks Add}

A network representation captures what traditional models cannot:
\begin{enumerate}
    \item \textbf{Direct linkages}: who owes money to whom, and how much.
    \item \textbf{Concentration risk}: if many institutions are exposed to the
          same counterparty, a single default can trigger widespread losses.
    \item \textbf{Cascade dynamics}: a default at node $i$ reduces the assets of
          node $i$'s creditors, potentially causing them to default, which in turn
          affects \emph{their} creditors.
    \item \textbf{Topology matters}: a highly connected ``hub'' bank failing is
          far more dangerous than a peripheral institution failing.
          Network centrality measures quantify this.
\end{enumerate}

\begin{keyidea}
Credit risk is not just about individual balance sheets.
It is about the \textbf{structure of obligations} between institutions.
Two banking systems with identical individual bank characteristics but different
network topologies can have dramatically different systemic risk profiles.
A network with a single highly connected hub is more fragile than a diversified
network with many moderate connections.
\end{keyidea}

\subsection{The Contribution of This Project}

We combine two powerful ideas:
\begin{enumerate}
    \item \textbf{Financial network models} (Eisenberg--Noe, DebtRank) that simulate
          how shocks cascade through interbank exposures.
    \item \textbf{Graph Neural Networks} that learn to predict which institutions
          will default by aggregating information across the network topology.
\end{enumerate}

The GNN approach is attractive because it can learn \emph{non-linear} contagion
patterns that simple cascade models may miss, and it can generalize to networks
with different topologies once trained.

%% ============================================================================
\section{Financial Networks}
\label{sec:financial-networks}
%% ============================================================================

A financial network is a directed, weighted graph $G = (V, E, W)$ where:
\begin{itemize}
    \item $V = \{1, 2, \ldots, n\}$ is the set of \textbf{nodes} (institutions:
          banks, insurance companies, hedge funds, corporates).
    \item $E \subseteq V \times V$ is the set of \textbf{directed edges} representing
          financial obligations.
    \item $W: E \to \RV_{>0}$ assigns a \textbf{weight} $w_{ij}$ to each edge
          $(i, j)$, representing the nominal value of $i$'s obligation to $j$
          (equivalently, $j$'s exposure to $i$).
\end{itemize}

\subsection{Types of Financial Linkages}

Multiple types of edges can coexist in a financial network:

\begin{center}
\begin{tabular}{@{}p{3.5cm}p{4cm}p{5cm}@{}}
\toprule
\textbf{Linkage Type} & \textbf{Edge Meaning} & \textbf{Example} \\
\midrule
Interbank lending     & Bank $i$ has borrowed from bank $j$
                      & Overnight lending, term deposits \\[4pt]
Derivative exposures  & Bank $i$ owes mark-to-market to $j$
                      & Interest rate swaps, CDS contracts \\[4pt]
Payment systems       & Bank $i$ sends payments through $j$
                      & Fedwire, TARGET2 \\[4pt]
Supply chain          & Firm $i$ supplies goods to firm $j$
                      & Trade credit, accounts receivable \\[4pt]
Common asset holdings & $i$ and $j$ hold overlapping portfolios
                      & Indirect link via fire-sale channel \\
\bottomrule
\end{tabular}
\end{center}

\subsection{The Liability Matrix}

The network is compactly represented by an $n \times n$ \textbf{liability matrix}
$L$, where $L_{ij}$ is the total obligation of institution $i$ to institution $j$:
\begin{equation}\label{eq:liability-matrix}
    L = \begin{pmatrix}
        0      & L_{12} & \cdots & L_{1n} \\
        L_{21} & 0      & \cdots & L_{2n} \\
        \vdots & \vdots & \ddots & \vdots \\
        L_{n1} & L_{n2} & \cdots & 0
    \end{pmatrix}
\end{equation}
The diagonal is zero (no self-obligations).
The total interbank liabilities of node $i$ are $\bar{p}_i = \sum_{j=1}^{n} L_{ij}$,
and the total interbank assets of node $i$ are $\sum_{j=1}^{n} L_{ji}$.

Each institution also has \textbf{external assets} $e_i$ (assets not related to
interbank claims) and \textbf{external liabilities} $d_i$ (deposits, bonds,
obligations to non-bank creditors).

\subsection{Topological Metrics}

Network topology determines how shocks propagate.
The following centrality measures quantify the importance of each node:

\begin{definition}[Centrality Measures]\label{def:centrality}
For a weighted directed graph with adjacency matrix $A$ (where $A_{ij} = w_{ij}$
if edge $(i,j)$ exists, 0 otherwise):
\begin{enumerate}
    \item \textbf{Degree centrality}: the number of connections.
          In-degree $k_i^{\text{in}} = \sum_j A_{ji}$ counts incoming edges
          (exposures \emph{to} $i$).
          Out-degree $k_i^{\text{out}} = \sum_j A_{ij}$ counts outgoing edges
          ($i$'s obligations).
    \item \textbf{Betweenness centrality}: the fraction of shortest paths between
          all pairs of nodes that pass through node $i$.
          A high-betweenness node is a critical intermediary; its failure disrupts
          many channels.
    \item \textbf{Eigenvector centrality}: solves $Ax = \lambda x$.
          A node is important if it is connected to other important nodes.
          The eigenvector corresponding to the largest eigenvalue assigns a centrality
          score to each node.
\end{enumerate}
\end{definition}

\begin{example}[Core-periphery structure]\label{ex:core-periphery}
Empirical interbank networks consistently exhibit \textbf{core-periphery structure}:

\begin{itemize}
    \item A small \textbf{core} of 10--15 major banks that are densely connected
          to each other (high degree, high betweenness).
    \item A large \textbf{periphery} of smaller institutions that connect to one or
          two core banks but not to each other.
\end{itemize}

This structure has important implications for contagion:
\begin{itemize}
    \item A core bank failing can propagate shocks to many peripheral banks
          simultaneously.
    \item A peripheral bank failing typically affects only one or two core banks,
          which can usually absorb the loss.
    \item The system is robust to random peripheral failures but fragile to
          targeted core failures --- the ``robust-yet-fragile'' property
          (Acemoglu, Ozdaglar \& Tahbaz-Salehi, 2015).
\end{itemize}
\end{example}

\subsection{Node Features}

Beyond topology, each node carries \textbf{features} that determine its resilience
to shocks.
For a GNN, these form the initial node feature matrix $X \in \RV^{n \times d}$
where $d$ is the number of features per node:

\begin{center}
\begin{tabular}{@{}lll@{}}
\toprule
\textbf{Feature} & \textbf{Symbol} & \textbf{Interpretation} \\
\midrule
Capital ratio       & $x_1$ & Equity / total assets; buffer against losses \\
Leverage            & $x_2$ & Total assets / equity; amplifies shocks \\
Asset quality       & $x_3$ & Non-performing loans / total loans \\
Size (log assets)   & $x_4$ & Systemic importance proxy \\
Liquidity ratio     & $x_5$ & Liquid assets / short-term liabilities \\
\bottomrule
\end{tabular}
\end{center}

A bank with a high capital ratio and low leverage can absorb larger losses before
defaulting.
The GNN will learn to combine these features with the network structure.

\begin{keyresult}
A financial network is fully described by three components:
\begin{enumerate}
    \item \textbf{Topology} --- the adjacency/liability matrix $L$: who owes whom,
          how much.
    \item \textbf{Node features} --- the feature matrix $X$: balance sheet
          characteristics of each institution.
    \item \textbf{External variables} --- external assets $e_i$ and liabilities $d_i$
          that anchor each institution's solvency outside the network.
\end{enumerate}
The GNN operates on all three simultaneously.
\end{keyresult}

%% ============================================================================
\section{Contagion Mechanisms}
\label{sec:contagion}
%% ============================================================================

Financial contagion operates through several distinct channels.
Understanding these channels is essential for designing the cascade simulation
that generates training labels for the GNN.

\subsection{Direct (Credit) Contagion}

The most transparent mechanism: institution $A$ defaults on its obligations,
and institution $B$, which is owed money by $A$, suffers a loss equal to its
exposure (net of recovery).

Formally, if institution $i$ defaults and $j$ has an exposure $L_{ij}$ to $i$
with recovery rate $R$, then $j$'s loss is:
\begin{equation}\label{eq:direct-loss}
    \ell_j = (1 - R) \cdot L_{ij}
\end{equation}
If $\ell_j$ exceeds $j$'s capital buffer $c_j$, then $j$ defaults as well,
and the cascade continues.

\begin{keyidea}
Direct contagion is a \textbf{domino effect}.
The key insight from Eisenberg \& Noe (2001) is that we cannot simply iterate
default-by-default, because obligations are \emph{simultaneous}: when $i$ defaults,
it may only partially pay its creditors, and the partial payment depends on how
much $i$ receives from \emph{its} debtors, who may themselves be defaulting.
This circular dependency requires solving for a \textbf{clearing payment vector}
--- a fixed-point problem.
\end{keyidea}

\subsubsection{The Eisenberg--Noe Clearing Mechanism}

Consider $n$ institutions with interbank liability matrix $L$.
Let $\bar{p}_i = \sum_{j} L_{ij}$ be the total interbank obligations of node~$i$,
and define the \textbf{relative liability matrix}:
\begin{equation}\label{eq:relative-liability}
    \Pi_{ij} = \frac{L_{ij}}{\bar{p}_i} \quad \text{for } \bar{p}_i > 0
\end{equation}
so that $\Pi_{ij}$ is the fraction of $i$'s total interbank obligations owed to $j$.

Each institution $i$ has external assets $e_i$ (cash, securities, loans to the
real economy).
The total resources available to $i$ are its external assets plus whatever it
receives from other institutions:
\begin{equation}\label{eq:total-resources}
    a_i(p) = e_i + \sum_{j=1}^{n} \Pi_{ji}\, p_j
\end{equation}
where $p_j$ is the actual payment made by institution $j$.

\begin{definition}[Clearing Payment Vector]\label{def:clearing}
A payment vector $p^* = (p_1^*, \ldots, p_n^*)$ is a \textbf{clearing vector} if
for each institution $i$:
\begin{equation}\label{eq:clearing}
    p_i^* = \min\!\bigl(\bar{p}_i,\; a_i(p^*)\bigr)
    = \min\!\left(\bar{p}_i,\; e_i + \sum_{j=1}^{n} \Pi_{ji}\, p_j^*\right)
\end{equation}
That is, each institution pays the minimum of what it owes ($\bar{p}_i$) and what
it has ($a_i$).
If $a_i(p^*) < \bar{p}_i$, institution $i$ is in \textbf{default} and pays pro
rata to its creditors.
\end{definition}

\begin{theorem}[Eisenberg \& Noe, 2001]\label{thm:eisenberg-noe}
Under mild regularity conditions (the financial system has no cycles of institutions
with zero external assets), the clearing payment vector $p^*$ exists and is unique.
It can be computed by the following \textbf{fixed-point iteration} (``fictitious
default algorithm''):

\begin{enumerate}
    \item Initialize $p^{(0)} = \bar{p}$ (every institution pays in full).
    \item At each iteration $k$, compute:
    \[
        p_i^{(k+1)} = \min\!\left(\bar{p}_i,\; e_i + \sum_{j=1}^{n} \Pi_{ji}\, p_j^{(k)}\right)
    \]
    \item Repeat until convergence: $\|p^{(k+1)} - p^{(k)}\| < \varepsilon$.
\end{enumerate}

The sequence $p^{(0)} \geq p^{(1)} \geq \cdots$ is non-increasing and converges
to $p^*$ in at most $n$ iterations (each iteration identifies at least one new
defaulting institution).
\end{theorem}

\begin{example}[Three-bank clearing]\label{ex:three-bank}
Three banks with liabilities:
\[
    L = \begin{pmatrix} 0 & 80 & 0 \\ 0 & 0 & 60 \\ 40 & 0 & 0 \end{pmatrix}
\]
So bank~1 owes 80 to bank~2, bank~2 owes 60 to bank~3, bank~3 owes 40 to bank~1.
External assets: $e_1 = 100$, $e_2 = 50$, $e_3 = 30$.

\textbf{No shock}: total assets of bank~1 = $e_1 + $ received from bank~3 = $100 + 40 = 140 > 80$
(pays in full).
Bank~2 = $50 + 80 = 130 > 60$ (pays in full).
Bank~3 = $30 + 60 = 90 > 40$ (pays in full).
No defaults.

\textbf{Shock}: reduce $e_1$ to 20 (asset loss of 80).
Now bank~1's total assets = $20 + 40 = 60 < 80$ (default; pays 60 pro rata to bank~2).
Bank~2 receives 60 instead of 80; total assets = $50 + 60 = 110 > 60$ (survives).
Bank~3 receives 60 from bank~2; total assets = $30 + 60 = 90 > 40$ (survives).
\textbf{One default}, no cascade.

\textbf{Larger shock}: reduce $e_1$ to 0.
Bank~1's total assets = $0 + 40 = 40 < 80$ (default; pays 40).
Bank~2 receives 40; total assets = $50 + 40 = 90 > 60$ (survives).
Still only one default, because bank~2 is well-capitalized.
\end{example}

\subsection{Fire-Sale Contagion}

When a distressed institution is forced to sell illiquid assets quickly, the sales
depress market prices.
Other institutions holding the same assets suffer mark-to-market losses even though
they are not directly connected to the distressed institution.

Let $\Delta q_k$ be the quantity of asset $k$ sold by the distressed institution.
If the price impact function is linear, the new price is:
\begin{equation}\label{eq:fire-sale}
    p_k' = p_k\!\left(1 - \frac{\Delta q_k}{\text{ADV}_k}\right)
\end{equation}
where $\text{ADV}_k$ is the average daily volume of asset $k$.
All institutions holding asset $k$ suffer a loss proportional to their holdings.

\begin{warningbox}
Fire-sale contagion is particularly dangerous because it operates through
\textbf{indirect linkages}: two institutions that have no direct contractual
relationship can infect each other if they hold overlapping portfolios.
This channel was central to the 2008 crisis, when forced selling of
mortgage-backed securities caused prices to spiral downward, inflicting
losses on institutions far removed from the original subprime defaults.
\end{warningbox}

\subsection{Funding (Liquidity) Contagion}

When an institution is perceived as risky, its funding counterparties may:
\begin{itemize}
    \item Refuse to roll over short-term funding (commercial paper, repo).
    \item Demand higher collateral (margin calls).
    \item Withdraw deposits.
\end{itemize}

This creates a self-fulfilling prophecy: the institution's liquidity position
deteriorates, forcing asset sales (fire-sale channel) or outright default,
even if it is fundamentally solvent.
Northern Rock (2007) and Bear Stearns (2008) exemplify pure funding contagion.

\begin{keyresult}
The three contagion channels interact and amplify each other:
\begin{enumerate}
    \item \textbf{Direct contagion} causes credit losses at counterparties.
    \item Weakened counterparties are forced to sell assets (\textbf{fire-sale contagion}).
    \item Market participants observe distress and withdraw funding
          (\textbf{funding contagion}).
    \item Funding withdrawal forces more asset sales, depressing prices further.
\end{enumerate}
For our GNN implementation, we focus on direct contagion via the Eisenberg--Noe
mechanism, as it is the most natural fit for graph-structured data.
The GNN's message-passing layers can implicitly learn patterns that echo
fire-sale and funding dynamics through shared features.
\end{keyresult}

%% ============================================================================
\section{Graph Neural Networks: The Key Idea}
\label{sec:gnn-key-idea}
%% ============================================================================

A Graph Neural Network learns representations of nodes (or edges, or the entire
graph) by \textbf{iteratively aggregating information from neighboring nodes}.
At each layer, a node's representation is updated based on its own features
\emph{and} the features of its neighbors.

\subsection{Why GNNs Fit Contagion}

Consider why this architecture is a natural match for credit contagion:
\begin{itemize}
    \item The risk of institution $i$ depends on the health of institutions
          \emph{connected to $i$} (its direct counterparties).
    \item But it also depends on the health of institutions connected to $i$'s
          counterparties (indirect exposure, two hops away).
    \item And so on, for longer chains of exposure.
\end{itemize}

A GNN with $k$ layers captures exactly $k$-hop neighborhood information.
After one layer, each node ``knows'' about its direct neighbors.
After two layers, each node has incorporated information from two-hop neighbors.
This mirrors how contagion propagates: one round of defaults, then second-round
effects, then third-round effects.

\begin{keyidea}
The message-passing mechanism in a GNN is a \textbf{learned analogue of cascade
simulation}.
Where the Eisenberg--Noe model uses a fixed rule (``pay what you can''),
the GNN learns a parametric aggregation function from data.
This allows it to capture non-linear interactions, heterogeneous node responses,
and patterns that a simple threshold model may miss.
\end{keyidea}

\subsection{The Message-Passing Framework}

All GNN architectures follow the same abstract \textbf{message-passing} pattern.
For each node $i$ at layer $l$:

\begin{equation}\label{eq:message-passing}
    h_i^{(l+1)} = \phi\!\left(h_i^{(l)},\; \bigoplus_{j \in \mathcal{N}(i)} \psi\!\left(h_i^{(l)}, h_j^{(l)}, e_{ij}\right)\right)
\end{equation}

where:
\begin{itemize}
    \item $h_i^{(l)} \in \RV^{d_l}$ is the \textbf{embedding} (hidden representation) of
          node $i$ at layer $l$.
          At layer 0, $h_i^{(0)} = x_i$ (the raw node features).
    \item $\mathcal{N}(i)$ is the set of \textbf{neighbors} of node $i$.
    \item $\psi(\cdot)$ is the \textbf{message function}: it computes a ``message''
          from neighbor $j$ to node $i$, potentially using the edge weight $e_{ij}$.
    \item $\bigoplus$ is an \textbf{aggregation operator} (sum, mean, max) that
          combines messages from all neighbors into a single vector.
    \item $\phi(\cdot)$ is the \textbf{update function}: it combines the aggregated
          messages with the node's current embedding to produce the new embedding.
\end{itemize}

Different GNN architectures correspond to different choices of $\psi$, $\bigoplus$,
and $\phi$.

\subsection{Key Architectures}

\begin{definition}[GCN --- Kipf \& Welling, 2017]\label{def:gcn}
The Graph Convolutional Network uses a simple spectral-inspired aggregation:
\begin{equation}\label{eq:gcn-layer}
    H^{(l+1)} = \sigma\!\left(\tilde{D}^{-1/2}\, \tilde{A}\, \tilde{D}^{-1/2}\, H^{(l)}\, W^{(l)}\right)
\end{equation}
where $\tilde{A} = A + I$ (adjacency matrix with added self-loops),
$\tilde{D}$ is the degree matrix of $\tilde{A}$
($\tilde{D}_{ii} = \sum_j \tilde{A}_{ij}$),
$H^{(l)} \in \RV^{n \times d_l}$ is the node feature matrix at layer $l$,
$W^{(l)} \in \RV^{d_l \times d_{l+1}}$ is a learnable weight matrix,
and $\sigma$ is a non-linear activation function (typically $\ReLU$).
\end{definition}

\begin{definition}[GraphSAGE --- Hamilton et al., 2017]\label{def:graphsage}
GraphSAGE samples a fixed number of neighbors and learns an aggregation function:
\begin{align}
    h_{\mathcal{N}(i)}^{(l)} &= \text{AGGREGATE}\!\left(\{h_j^{(l)} : j \in \mathcal{N}(i)\}\right) \label{eq:sage-agg}\\
    h_i^{(l+1)} &= \sigma\!\left(W^{(l)} \cdot \bigl[h_i^{(l)} \,\|\, h_{\mathcal{N}(i)}^{(l)}\bigr]\right) \label{eq:sage-update}
\end{align}
where $\|$ denotes concatenation and AGGREGATE can be mean, LSTM, or max-pooling.
The concatenation of self-features with neighbor-features explicitly separates
the node's own information from its neighborhood signal.
\end{definition}

\begin{definition}[GAT --- Veli\v{c}kovi\'c et al., 2018]\label{def:gat}
The Graph Attention Network learns attention weights that determine how much
each neighbor contributes:
\begin{align}
    \alpha_{ij} &= \frac{\exp\!\bigl(\text{LeakyReLU}(a^T [W h_i \,\|\, W h_j])\bigr)}{\sum_{k \in \mathcal{N}(i)} \exp\!\bigl(\text{LeakyReLU}(a^T [W h_i \,\|\, W h_k])\bigr)} \label{eq:gat-attn}\\[4pt]
    h_i^{(l+1)} &= \sigma\!\left(\sum_{j \in \mathcal{N}(i)} \alpha_{ij}\, W\, h_j^{(l)}\right) \label{eq:gat-update}
\end{align}
where $a$ is a learnable attention vector.
The attention weights $\alpha_{ij}$ allow the model to focus on the most
relevant neighbors --- for credit contagion, it might learn to pay more attention
to highly leveraged counterparties.
\end{definition}

\begin{remark}
For credit contagion, GAT is conceptually appealing because not all counterparties
are equally dangerous.
A highly leveraged neighbor is a greater contagion risk than a well-capitalized one.
Attention weights can learn this distinction from data without explicit feature
engineering.
\end{remark}

%% ============================================================================
\section{GCN for Default Prediction}
\label{sec:gcn-default}
%% ============================================================================

We now develop the specific GCN architecture for predicting which nodes in a
financial network will default following an exogenous shock.

\subsection{Problem Formulation}

\textbf{Given}:
\begin{itemize}
    \item A financial network $G = (V, E)$ with $n$ nodes.
    \item Adjacency matrix $A \in \RV^{n \times n}$ (derived from the liability matrix;
          $A_{ij} = 1$ if $L_{ij} > 0$, or $A_{ij} = L_{ij}$ for a weighted graph).
    \item Node feature matrix $X \in \RV^{n \times d}$ with $d$ features per node
          (capital ratio, leverage, asset quality, size, liquidity ratio).
    \item A shock scenario $s$ that reduces the external assets of one or more nodes.
\end{itemize}

\textbf{Predict}: for each node $i$, the probability $\hat{y}_i \in [0, 1]$ that
node $i$ defaults under shock scenario $s$.

This is a \textbf{node-level binary classification} task on a graph.

\subsection{Two-Layer GCN Architecture}

We use a two-layer GCN, which captures information from 2-hop neighborhoods:

\textbf{Layer 1} (feature extraction):
\begin{equation}\label{eq:gcn-layer1}
    H^{(1)} = \ReLU\!\left(\tilde{D}^{-1/2}\, \tilde{A}\, \tilde{D}^{-1/2}\, X\, W^{(0)}\right)
\end{equation}
This maps from the original $d = 5$ features to a hidden dimension $d_1 = 16$.
So $W^{(0)} \in \RV^{5 \times 16}$.

\textbf{Layer 2} (classification):
\begin{equation}\label{eq:gcn-layer2}
    Z = \sigmoid\!\left(\tilde{D}^{-1/2}\, \tilde{A}\, \tilde{D}^{-1/2}\, H^{(1)}\, W^{(1)}\right)
\end{equation}
This maps from the hidden dimension $d_1 = 16$ to a single output.
So $W^{(1)} \in \RV^{16 \times 1}$, and $Z \in \RV^{n \times 1}$ contains the
predicted default probabilities.

The sigmoid activation in the final layer ensures outputs are in $[0, 1]$.

\subsection{What the Normalization Does}

The symmetric normalization $\tilde{D}^{-1/2}\, \tilde{A}\, \tilde{D}^{-1/2}$
has an intuitive interpretation.
For each node $i$, it computes a weighted average of its neighbors' features
(including its own, via the self-loop in $\tilde{A}$):

\begin{equation}\label{eq:gcn-nodeview}
    h_i^{(l+1)} = \sigma\!\left(\sum_{j \in \mathcal{N}(i) \cup \{i\}}
    \frac{1}{\sqrt{\tilde{d}_i\, \tilde{d}_j}}\; h_j^{(l)}\; W^{(l)}\right)
\end{equation}

where $\tilde{d}_i = 1 + \sum_j A_{ij}$ is the degree of node $i$ in the
augmented graph.
The $1/\sqrt{\tilde{d}_i\, \tilde{d}_j}$ factor normalizes by the degrees of both
the sending and receiving nodes, preventing high-degree nodes from dominating
the aggregation.

\begin{example}[Worked example: 50-node network]\label{ex:50-node}
Consider a synthetic interbank network with $n = 50$ nodes and $d = 5$ features
per node.

\textbf{Network}: generated with a core-periphery structure.
10 core banks are densely connected (edge probability 0.7 within core).
40 peripheral banks each connect to 1--3 core banks.
Total edges: approximately 100.

\textbf{Features}: for each node, we generate $x_i \in \RV^5$:
capital ratio (mean 0.08, std 0.03), leverage (mean 15, std 5),
NPL ratio (mean 0.03, std 0.02), log-assets (mean 10, std 2),
liquidity ratio (mean 0.25, std 0.10).

\textbf{Architecture}:
\begin{itemize}
    \item Input: $X \in \RV^{50 \times 5}$.
    \item Layer 1: $H^{(1)} = \ReLU(\hat{A}\, X\, W^{(0)}) \in \RV^{50 \times 16}$,
          where $\hat{A} = \tilde{D}^{-1/2}\tilde{A}\tilde{D}^{-1/2}$.
    \item Layer 2: $Z = \sigmoid(\hat{A}\, H^{(1)}\, W^{(1)}) \in \RV^{50 \times 1}$.
\end{itemize}

\textbf{Parameters}: $W^{(0)}$ has $5 \times 16 = 80$ entries, $W^{(1)}$ has
$16 \times 1 = 16$ entries.
Total learnable parameters: \textbf{96}.
This is remarkably compact --- the network structure provides the inductive bias
that keeps the model small.

\textbf{Training}: we generate 1{,}000 shock scenarios using the Eisenberg--Noe
cascade simulation (Section~\ref{sec:cascade}), producing binary default labels
for each node under each scenario.
Train with binary cross-entropy loss and the Adam optimizer.
\end{example}

\subsection{Loss Function}

The binary cross-entropy loss for node-level classification:
\begin{equation}\label{eq:bce}
    \mathcal{L} = -\frac{1}{n}\sum_{i=1}^{n}\Bigl[y_i \log \hat{y}_i + (1 - y_i)\log(1 - \hat{y}_i)\Bigr]
\end{equation}
where $y_i \in \{0, 1\}$ is the true default label (from the cascade simulation)
and $\hat{y}_i = Z_i$ is the GCN's predicted default probability.

\begin{warningbox}
Default events are typically \textbf{rare}: in a 50-node network, a typical shock
might cause 2--5 defaults.
This means the dataset is heavily imbalanced (90\%+ non-default labels).
Standard remedies apply:
\begin{itemize}
    \item \textbf{Weighted loss}: increase the weight on the positive (default) class.
    \item \textbf{Oversampling}: generate more severe shock scenarios to increase
          the default rate.
    \item \textbf{Focal loss}: down-weight easy (non-default) examples.
\end{itemize}
Without addressing class imbalance, the GCN may learn to predict ``no default''
for every node and achieve high accuracy but zero recall.
\end{warningbox}

%% ============================================================================
\section{Cascade Simulation}
\label{sec:cascade}
%% ============================================================================

The cascade simulation generates the \textbf{training data} for the GNN\@.
We simulate the Eisenberg--Noe clearing mechanism under many shock scenarios
to produce binary default labels for each node.

\subsection{Simulation Procedure}

\begin{enumerate}
    \item \textbf{Generate network}: create an $n$-node financial network with
          realistic core-periphery topology.
          Assign edge weights (interbank liabilities) and node features
          (capital ratio, leverage, etc.).
          Assign external assets $e_i$ to each node.

    \item \textbf{Sample shock}: select one or more ``seed'' nodes to shock.
          Reduce their external assets by a fraction $\delta \in [0, 1]$:
          \begin{equation}\label{eq:shock}
              e_i^{\text{shocked}} = (1 - \delta)\, e_i
          \end{equation}
          Vary $\delta$ and the choice of seed nodes across scenarios.

    \item \textbf{Run clearing}: apply the Eisenberg--Noe fixed-point iteration
          (Theorem~\ref{thm:eisenberg-noe}) to find the clearing payment vector $p^*$.

    \item \textbf{Identify defaults}: institution $i$ is in default if
          $p_i^* < \bar{p}_i$ (it cannot pay its obligations in full).
          Assign label $y_i = 1$ (default) or $y_i = 0$ (solvent).

    \item \textbf{Record}: store the tuple $(X, A, y)$ where $X$ is the
          (potentially shocked) node feature matrix, $A$ is the adjacency matrix,
          and $y \in \{0, 1\}^n$ is the vector of default labels.

    \item \textbf{Repeat}: generate $M$ scenarios (e.g., $M = 1{,}000$) by varying
          the shock location and severity.
\end{enumerate}

\subsection{Shock Design}

The choice of shocks determines the diversity and difficulty of the training data:

\begin{center}
\begin{tabular}{@{}p{3.5cm}p{4cm}p{5cm}@{}}
\toprule
\textbf{Shock Type} & \textbf{Description} & \textbf{Purpose} \\
\midrule
Single-node         & Shock one randomly chosen node
                    & Tests direct contagion from each node \\[4pt]
Multi-node          & Shock 2--5 randomly chosen nodes
                    & Tests overlapping cascades \\[4pt]
Targeted core       & Shock a high-centrality core node
                    & Tests worst-case systemic scenarios \\[4pt]
Severity sweep      & Fix the shocked node, vary $\delta$ from 0.2 to 1.0
                    & Maps the contagion threshold \\
\bottomrule
\end{tabular}
\end{center}

\begin{example}[Cascade statistics]\label{ex:cascade-stats}
For a 50-node core-periphery network with 10 core banks and 40 peripheral banks:

\begin{center}
\begin{tabular}{@{}lcc@{}}
\toprule
\textbf{Shock scenario} & \textbf{Avg.\ defaults} & \textbf{Max defaults} \\
\midrule
Single peripheral, $\delta = 0.5$  & 1.2 & 3  \\
Single peripheral, $\delta = 1.0$  & 1.8 & 5  \\
Single core, $\delta = 0.5$        & 3.4 & 8  \\
Single core, $\delta = 1.0$        & 6.1 & 15 \\
Two core banks, $\delta = 0.8$     & 9.3 & 22 \\
\bottomrule
\end{tabular}
\end{center}

\textbf{Observation}: shocking a core bank produces far more cascading defaults
than shocking a peripheral bank at the same severity.
The network topology amplifies shocks that hit the core.
\end{example}

\begin{keyresult}
The cascade simulation serves two purposes:
\begin{enumerate}
    \item \textbf{Training data}: the $(X, A, y)$ tuples are used to train the GCN
          to predict defaults from node features and network topology.
    \item \textbf{Benchmark}: the clearing mechanism provides ground-truth default
          labels against which the GCN's predictions are evaluated.
          If the GCN can learn to predict cascade outcomes accurately, it can be
          used as a \textbf{fast surrogate} for the computationally expensive clearing
          simulation --- enabling real-time systemic risk assessment.
\end{enumerate}
\end{keyresult}

%% ============================================================================
\section{Systemic Risk Metrics}
\label{sec:systemic-metrics}
%% ============================================================================

Once we can simulate cascades (or predict them with a trained GNN), we can compute
metrics that quantify the systemic importance of each institution and the overall
fragility of the network.

\subsection{DebtRank}

DebtRank (Battiston et al., 2012) measures the fraction of total economic value
in the network that is lost when a given institution enters distress.

\begin{definition}[DebtRank]\label{def:debtrank}
Let $v_i$ denote the total economic value (equity + interbank assets) of institution~$i$,
and let $V = \sum_i v_i$ be the total system value.
When institution $i$ is shocked into distress, the cascade propagates and
causes losses at other institutions.
Let $L_i^{\text{total}}$ be the total loss across the system (including at $i$
itself).
The DebtRank of institution $i$ is:
\begin{equation}\label{eq:debtrank}
    \DebtRank(i) = \frac{L_i^{\text{total}} - v_i}{V - v_i}
\end{equation}
This is the fraction of \textbf{other institutions'} value lost due to the distress
of $i$.
A DebtRank of 0 means $i$'s failure has no systemic impact.
A DebtRank of 1 means $i$'s failure wipes out the entire rest of the system.
\end{definition}

The DebtRank computation uses a \emph{linear contagion} model: when institution $j$
is exposed to distressed institution $i$, $j$'s equity is reduced proportionally
to the exposure:
\begin{equation}\label{eq:debtrank-prop}
    h_j^{(t+1)} = \min\!\left(1,\; h_j^{(t)} + \sum_{i:\, s_i^{(t)} = \text{active}} W_{ji}\, h_i^{(t)}\right)
\end{equation}
where $h_i^{(t)} \in [0, 1]$ is the ``distress level'' of institution $i$ at
round $t$, $W_{ji} = L_{ij} / v_j$ is the relative exposure of $j$ to $i$,
and $s_i^{(t)}$ tracks whether $i$ is still actively propagating distress.

\begin{remark}
DebtRank differs from the Eisenberg--Noe model in a crucial way: it uses
\textbf{continuous distress levels} $h_i \in [0, 1]$ rather than binary default.
This captures the economic damage from a bank's stock price declining by 50\%
(partial distress) even if it does not technically default.
It is therefore a more sensitive measure of systemic impact.
\end{remark}

\subsection{CoVaR}

CoVaR (Adrian \& Brunnermeier, 2016) measures the tail risk of the financial system
conditional on the distress of a specific institution.

\begin{definition}[CoVaR]\label{def:covar}
The $\CoVaR_q^{j|i}$ is the $q$-quantile of the loss distribution of institution
(or system) $j$, \textbf{conditional} on institution $i$ being in distress
(typically defined as $i$'s return being at its $\VaR_q$ level):
\begin{equation}\label{eq:covar}
    \PV\!\left(X_j \leq \CoVaR_q^{j|i} \;\Big|\; X_i = \VaR_q^i\right) = q
\end{equation}
The systemic risk contribution of institution $i$ is:
\begin{equation}\label{eq:delta-covar}
    \Delta\CoVaR_q^{j|i} = \CoVaR_q^{j|i} - \CoVaR_q^{j|\text{median}(i)}
\end{equation}
This measures how much worse the system's tail risk becomes when institution $i$
is in distress compared to when $i$ is at its median state.
\end{definition}

In our network context, $X_j$ can be the total system loss (sum of losses across
all institutions), and we condition on institution $i$ being at or beyond its
$\VaR$ threshold.
The cascade simulation naturally produces these conditional distributions:
run the cascade with $i$ shocked, record system-wide losses, and compute quantiles.

\subsection{SRISK}

SRISK measures the expected capital shortfall of an institution in a systemic crisis.

\begin{definition}[SRISK]\label{def:srisk}
\begin{equation}\label{eq:srisk}
    \SRISK_i = \EV\!\bigl[\max\bigl(k \cdot A_i - E_i,\; 0\bigr) \;\big|\; \text{crisis}\bigr]
\end{equation}
where $k$ is a prudential capital ratio (e.g., 8\%), $A_i$ is the total assets of
institution $i$ under the crisis scenario, and $E_i$ is its equity.
$\SRISK_i > 0$ means institution $i$ would need a capital injection in a crisis.
The total system SRISK, $\sum_i \SRISK_i$, is a measure of the system's overall
capital shortfall and hence its vulnerability.
\end{definition}

\subsection{Using the GNN to Compute Systemic Metrics}

Once the GCN is trained, it can serve as a fast proxy for the cascade simulation:

\begin{enumerate}
    \item \textbf{DebtRank via GNN}: for each node $i$, create a shock scenario
          (reduce $i$'s capital), run the GCN to predict default probabilities
          for all other nodes, and compute the expected system-wide loss.
          This is orders of magnitude faster than running the full
          Eisenberg--Noe clearing for each node.

    \item \textbf{CoVaR via GNN}: generate many shock scenarios, use the GCN to
          predict default outcomes, and compute the conditional VaR of system
          losses given distress at node $i$.

    \item \textbf{Stress testing}: apply a macroeconomic shock scenario (e.g.,
          all banks' external assets decrease by 20\%), feed the shocked features
          into the GCN, and obtain a real-time map of which institutions are
          most likely to default.
\end{enumerate}

\begin{keyidea}
The GNN is not just a prediction model --- it is a \textbf{computational accelerator}
for systemic risk analysis.
Full cascade simulations on large networks (thousands of nodes) are expensive.
A trained GNN performs the same prediction in a single forward pass (two matrix
multiplications), enabling real-time scenario analysis, interactive stress testing,
and optimization of network structure for resilience.
\end{keyidea}

\begin{refbox}
\begin{itemize}
    \item For a rigorous treatment of DebtRank and its extensions, see
          Bardoscia et al.\ (2017), ``Pathways towards instability in financial
          networks,'' \textit{Nature Communications} 8, 14416.
    \item For the connection between GNNs and financial stability analysis, see
          the survey by Xu et al.\ (2021), ``Graph Neural Networks for Financial
          Network Analysis.''
    \item For the theoretical foundations of network contagion, see
          Acemoglu, Ozdaglar \& Tahbaz-Salehi (2015), ``Systemic Risk and
          Stability in Financial Networks,'' \textit{American Economic Review}
          105(2), 564--608.
\end{itemize}
\end{refbox}

%% ============================================================================
\section{References}
\label{sec:references}
%% ============================================================================

The following works provide the theoretical foundations for the models and methods
developed in this document.

\begin{enumerate}[label={[\arabic*]}, leftmargin=2cm]
    \item \textbf{Eisenberg, L.\ \& Noe, T.H.}\ (2001).
          ``Systemic Risk in Financial Systems.''
          \emph{Management Science}, 47(2), 236--249.
          The foundational clearing mechanism for interbank payment systems.
          Proves existence and uniqueness of the clearing payment vector and establishes
          the fixed-point algorithm used in Section~\ref{sec:cascade}.

    \item \textbf{Battiston, S., Puliga, M., Kaushik, R., Tasca, P.\ \& Caldarelli, G.}\ (2012).
          ``DebtRank: Too Central to Fail? Financial Networks, the FED and Systemic Risk.''
          \emph{Scientific Reports}, 2, 541.
          Introduces the DebtRank metric and applies it to the network of emergency
          lending during the 2008 crisis, showing that the largest banks also had
          the highest systemic impact.

    \item \textbf{Kipf, T.N.\ \& Welling, M.}\ (2017).
          ``Semi-Supervised Classification with Graph Convolutional Networks.''
          \emph{Proceedings of ICLR 2017}.
          The seminal GCN paper.
          Derives the spectral motivation for the $\tilde{D}^{-1/2}\tilde{A}\tilde{D}^{-1/2}$
          normalization and demonstrates its effectiveness on citation networks.

    \item \textbf{Adrian, T.\ \& Brunnermeier, M.K.}\ (2016).
          ``CoVaR.''
          \emph{American Economic Review}, 106(7), 1705--1741.
          Defines CoVaR and $\Delta$CoVaR as systemic risk measures and estimates
          them for major financial institutions using quantile regression.

    \item \textbf{Acemoglu, D., Ozdaglar, A.\ \& Tahbaz-Salehi, A.}\ (2015).
          ``Systemic Risk and Stability in Financial Networks.''
          \emph{American Economic Review}, 105(2), 564--608.
          Shows that the relationship between connectivity and stability is
          non-monotone: moderate diversification enhances stability, but beyond
          a threshold, more connections propagate shocks more widely.

    \item \textbf{Hamilton, W.L., Ying, R.\ \& Leskovec, J.}\ (2017).
          ``Inductive Representation Learning on Large Graphs.''
          \emph{Proceedings of NeurIPS 2017}.
          Introduces GraphSAGE, a scalable GNN framework that samples and
          aggregates neighborhood features.

    \item \textbf{Veli\v{c}kovi\'c, P., Cucurull, G., Casanova, A., Romero, A.,
          Li\`o, P.\ \& Bengio, Y.}\ (2018).
          ``Graph Attention Networks.''
          \emph{Proceedings of ICLR 2018}.
          Introduces attention-based aggregation for GNNs, enabling learned
          weighting of neighbor contributions.

    \item \textbf{Bardoscia, M., Battiston, S., Caccioli, F.\ \& Caldarelli, G.}\ (2017).
          ``Pathways towards instability in financial networks.''
          \emph{Nature Communications}, 8, 14416.
          Extends DebtRank to incorporate multiple contagion channels and shows
          how small shocks can trigger large cascades through non-linear amplification.
\end{enumerate}

\begin{prereq}
Before implementing the GCN for credit contagion in code, ensure you are comfortable with:
\begin{itemize}
    \item Basic graph theory: adjacency matrices, degree, paths, connected components.
    \item Neural network fundamentals: forward pass, backpropagation, loss functions,
          gradient descent.
    \item PyTorch Geometric (or DGL): constructing graph data objects, defining GNN
          layers, training loops.
    \item Credit risk basics: default, recovery rate, capital adequacy, leverage.
    \item Linear algebra: matrix multiplication, eigenvalues (for the spectral motivation
          behind GCN).
\end{itemize}
\end{prereq}

\vfill
\hrule
\smallskip
\noindent\textit{Andr\'es Gonz\'alez Ortega --- Risk Analyst Project} \hfill \textit{March 2026}

\end{document}
